\documentclass[a4paper, serif, 10pt]{moderncv}

%% moderncv
% color and style
\moderncvcolor{black}
\moderncvstyle{banking}

%% page layout/margins
\usepackage[top=2.5cm, bottom=2.5cm, left=1.5cm, right=1.5cm]{geometry}

\nopagenumbers{}

%% Babel config for Spanish language
% \usepackage[spanish]{babel} has some conflicting issues with moderncv
% so we have to use enable the following options to make it work
%
% ref:
% - https://tex.stackexchange.com/a/140161/36007
\usepackage[spanish,es-lcroman]{babel}

%% fontspec
\usepackage{fontspec}

\IfFontExistsTF{Linux Libertine}{
  \setmainfont[Ligatures={TeX, Common}, Numbers=OldStyle]{Linux Libertine}
}{}
\IfFontExistsTF{Linux Libertine O}{
  \setmainfont[Ligatures={TeX, Common}, Numbers=OldStyle]{Linux Libertine O}
}{}

%% CTeX
% CJK support, used to show CJK characters in the resume
%
% - heading=false: disable ctex heading
% - punct=kaiming: use kaiming punctuations styles for CJK
% - scheme=plain: use plain scheme, do not override \normalsize font size
% - space=auto: space settings for CJK characters
%
% CJK fonts are resolved per locale language, e.g. Noto Serif/Sans CJK JP
% for Japanese resumes, with fallback to the next region if not installed
%
% ref:
% - http://ctan.mirrorcatalogs.com/language/chinese/ctex/ctex.pdf
\usepackage[UTF8, heading=false, punct=kaiming, scheme=plain, space=auto]{ctex}

\IfFontExistsTF{Noto Serif CJK SC}{
  \setCJKmainfont{Noto Serif CJK SC}
}{}
\IfFontExistsTF{Noto Sans CJK SC}{
  \setCJKsansfont{Noto Sans CJK SC}
}{}

%% line spacing
\usepackage{setspace}
\setstretch{1.125}

%% URL styling - use normal text instead of monospace
\urlstyle{same}

%% Patch \httplink and \httpslink to handle full URLs
% The original moderncv macros always prepend http:// or https:// to URLs.
% This causes issues when the URL already has a protocol (e.g., https://example.com)
% resulting in malformed URLs like https://https://example.com.
% This patch checks if the URL already contains "://" and uses it directly if so.
% Note: We use \str_set:Nx (x-expansion) to expand macros like \@homepage first.
\ExplSyntaxOn
\str_new:N \g_yamlresume_url_str

\RenewDocumentCommand{\httpslink}{O{}m}{
  \str_set:Nx \g_yamlresume_url_str {#2}
  \str_if_in:NnTF \g_yamlresume_url_str {://}
    { \url{#2} }
    { \url{https://#2} }
}

\RenewDocumentCommand{\httplink}{O{}m}{
  \str_set:Nx \g_yamlresume_url_str {#2}
  \str_if_in:NnTF \g_yamlresume_url_str {://}
    { \url{#2} }
    { \url{http://#2} }
}
\ExplSyntaxOff

\name{Lucía Fernández Moreno}{}
\title{Gerente de Operaciones | Eficiencia, personas y mejora continua}
\phone[mobile]{+34 612 345 678}
\email{lucia.fernandez@example.com}
\homepage{https://www.luciafernandez.es}

\address{Calle de Alcalá 120, 4º B, Madrid, Comunidad de Madrid, España, 28009}{}{}

\extrainfo{{\small \faLinkedin}\ \href{https://www.linkedin.com/in/lucia-fernandez-ops}{@lucia-fernandez-ops} {} {} {} • {} {} {} 
{\small \faMedium}\ \href{https://medium.com/@luciafernandezops}{@luciafernandezops}}

\begin{document}

\maketitle

\section{Información básica}

\cvline{}{\begin{itemize}
\item Gerente de Operaciones con más de 10 años de experiencia en logística, retail y servicios, liderando equipos multidisciplinares de hasta 45 personas.
\item Especializada en optimización de procesos, control de costes, planificación de la demanda y mejora continua mediante metodologías Lean Six Sigma.
\item Historial demostrado en la reducción de tiempos de entrega, aumento de la productividad y despliegue de indicadores operativos alineados con objetivos estratégicos.
\item Profesional orientada a resultados, con sólida capacidad de negociación, gestión de proveedores y coordinación entre áreas comerciales, financieras y de atención al cliente.
\end{itemize}}
\section{Educación}

\cventry{sept 2013–jul 2014}
        {Maestría, Dirección de Operaciones, Puntuación: 8,9/10}
        {Universidad Carlos III de Madrid}
        {\url{https://www.uc3m.es}}
        {}
        {\begin{itemize}
\item Formación avanzada en diseño, planificación y control de operaciones.
\item Trabajo de fin de máster sobre optimización de inventarios en retail.
\end{itemize}
\textbf{Cursos}: Gestión de la Cadena de Suministro, Dirección de Producción, Análisis de Datos para Operaciones, Gestión de Calidad, Negociación y Gestión de Proveedores}

\cventry{sept 2009–jun 2013}
        {Licenciatura, Administración y Dirección de Empresas, Puntuación: 8,2/10}
        {Universidad de Valladolid}
        {\url{https://www.uva.es}}
        {}
        {\begin{itemize}
\item Base sólida en gestión empresarial, finanzas y análisis cuantitativo.
\item Participación en proyectos de consultoría para pymes locales.
\end{itemize}
\textbf{Cursos}: Economía de la Empresa, Contabilidad Financiera, Estadística Aplicada, Dirección de Recursos Humanos, Logística Empresarial}
\section{Experiencia laboral}

\cventry{mar 2019 hasta la fecha}
        {Gerente de Operaciones}
        {Distribuciones NovaRetail S.L.}
        {\url{https://www.novaretail.es}}
        {}
        {\begin{itemize}
\item Dirección de las operaciones de una red de 18 tiendas y un centro de distribución, con un equipo de 45 personas.
\item Rediseño de los procesos de reposición y planificación de la demanda, reduciendo el stock obsoleto un 23 \% y mejorando el nivel de servicio al 97 \%.
\item Implantación de un cuadro de mando con indicadores de productividad, costes logísticos y satisfacción del cliente.
\item Negociación de contratos con proveedores y operadores logísticos, con un ahorro anual superior a 180.000 EUR.
\item Impulso de programas de formación y desarrollo del equipo, disminuyendo la rotación un 15 \%.
\end{itemize}
\textbf{Palabras clave}: Lean Six Sigma, Planificación S\&OP, Control de costes, Gestión de equipos, Indicadores}

\cventry{jun 2015–feb 2019}
        {Coordinadora de Operaciones}
        {Logística Ibérica Express}
        {\url{https://www.logistica-iberica.com}}
        {}
        {\begin{itemize}
\item Coordinación de una flota de 60 vehículos y de las operaciones diarias de reparto en la zona centro.
\item Optimización de rutas y planificación de turnos, logrando una reducción del 12 \% en kilómetros recorridos.
\item Seguimiento de acuerdos de nivel de servicio con clientes clave y resolución de incidencias operativas.
\item Colaboración con el área de sistemas en la digitalización de albaranes y pruebas de entrega.
\end{itemize}
\textbf{Palabras clave}: Gestión de flotas, Optimización de rutas, Atención al cliente, Nivel de servicio}

\cventry{sept 2014–may 2015}
        {Supervisora de Tienda}
        {Grupo Mercado Central}
        {\url{https://www.grupomercadocentral.es}}
        {}
        {\begin{itemize}
\item Supervisión de un equipo de 12 personas en una tienda de gran superficie.
\item Gestión de inventarios, reposición y control de mermas.
\item Formación en atención al cliente y cumplimiento de estándares de calidad y seguridad.
\end{itemize}
\textbf{Palabras clave}: Gestión de inventario, Atención al cliente, Control de mermas, Seguridad}
\section{Idiomas}

\cvline{Español}{Competencia nativa o bilingüe \hfill \textbf{Palabras clave}: Español nativo}
\cvline{Inglés}{Competencia profesional plena \hfill \textbf{Palabras clave}: C1, Uso profesional}
\cvline{Francés}{Competencia limitada de trabajo \hfill \textbf{Palabras clave}: A2}
\section{Competencias}

\cvline{Gestión de Operaciones}{Experto \hfill \textbf{Palabras clave}: Planificación de demanda, Mejora continua, Lean Six Sigma, Gestión de inventario, Planificación S\&OP}
\cvline{Liderazgo y Equipos}{Avanzado \hfill \textbf{Palabras clave}: Gestión de equipos, Mentoría, Negociación, Gestión del cambio, Comunicación}
\cvline{Análisis y Herramientas}{Avanzado \hfill \textbf{Palabras clave}: Excel avanzado, Power BI, SAP, Tableau, SQL}
\cvline{Logística y Cadena de Suministro}{Experto \hfill \textbf{Palabras clave}: Aprovisionamiento, Gestión de proveedores, Almacenamiento, Transporte, Última milla}
\section{Premios}

\cventry{nov 2022}
        {Distribuciones NovaRetail S.L.}
        {Premio a la Excelencia Operativa}
        {}
        {}
        {Reconocimiento por la mejora del nivel de servicio y la reducción de costes logísticos durante el ejercicio 2022.}

\cventry{jun 2018}
        {Logística Ibérica Express}
        {Reconocimiento por Mejora Continua}
        {}
        {}
        {Distinción por la optimización de rutas y la disminución de incidencias en la zona centro.}
\section{Certificados}

\cventry{feb 2020}
        {ASQ}
        {Lean Six Sigma Green Belt}
        {\url{https://asq.org/cert/lean-six-sigma-green-belt}}
        {}
        {}

\cventry{sept 2019}
        {APICS}
        {Certificación en Gestión de Cadena de Suministro}
        {\url{https://www.ascm.org/credentials/cscp/}}
        {}
        {}
\section{Publicaciones}

\cventry{abr 2021}
        {Optimización de inventarios en redes de retail}
        {Revista de Dirección y Operaciones}
        {\url{https://www.revistadireccionoperaciones.com/articulos/optimizacion-inventarios-retail}}
        {}
        {\begin{itemize}
\item Artículo sobre modelos de reposición y planificación de la demanda en entornos multitienda.
\item Análisis de indicadores clave para reducir roturas de stock sin aumentar el capital inmovilizado.
\end{itemize}}
\section{Proyectos}

\cventry{ene 2021–oct 2021}
        {Implantación de un sistema de inventario automatizado en el centro de distribución.}
        {Automatización del Inventario}
        {\url{https://www.novaretail.es/proyectos/inventario-automatizado}}
        {}
        {\begin{itemize}
\item Definición de requisitos junto con sistemas y operaciones.
\item Pruebas piloto, formación de usuarios y despliegue en 18 tiendas.
\item Reducción del tiempo de recuento en un 40 \% y mejora de la exactitud del stock.
\end{itemize}
\textbf{Palabras clave}: Automatización, Inventario, Formación, Procesos}

\cventry{mar 2018–dic 2018}
        {Programa para reducir mermas en tienda y almacén mediante análisis de causas y acciones correctivas.}
        {Programa de Reducción de Mermas}
        {\url{https://www.logistica-iberica.com/casos/mermas}}
        {}
        {\begin{itemize}
\item Análisis de datos de mermas por categoría y tienda.
\item Implantación de controles de recepción y formación al personal.
\item Disminución de la merma un 18 \% en nueve meses.
\end{itemize}
\textbf{Palabras clave}: Mejora continua, Control de mermas, Análisis de datos}
\section{Intereses}

\cvline{Deporte}{Carrera, Senderismo, Ciclismo}
\cvline{Lectura}{Novela histórica, Ensayo, Economía}
\section{Voluntariado}

\cventry{ene 2017–dic 2020}
        {Voluntaria de Logística}
        {Fundación Banco de Alimentos de Madrid}
        {\url{https://www.bamadrid.org}}
        {}
        {\begin{itemize}
\item Apoyo en la coordinación de la recepción y clasificación de alimentos.
\item Colaboración en la planificación de rutas de reparto a entidades sociales.
\item Participación en campañas de recogida y formación de nuevos voluntarios.
\end{itemize}}
    
\end{document}